%*******************************************************
% Summary
%*******************************************************
\pdfbookmark[1]{Zusammenfassung}{summary}
\chapter*{Zusammenfassung}
\addcontentsline{toc}{chapter}{Zusammenfassung}

Zusammenfassung der Arbeit.
Der Abschnitt dient dem Leser zur schnellen Orientierung und Einordnung der vorliegenden Arbeit.
Bei der Formulierung sollten wichtige Schlüsselbegriffe verwendet werden.

Checkliste:
\begin{enumerate}
\item Kurze, klare Darstellung zu Gegenstand, Fragestellung und Zielsetzung der Arbeit
\item Angewandte Methodik
\item Vorstellung der wesentlichen Ergebnisse und daraus resultierende Schlussfolgerungen.
\end{enumerate}
Für wissenschaftliche Abschlussarbeiten sind etwa 200 bis 250 Wörter ausreichend.

\paragraph{}
Idealerweise spart Ihnen diese Vorlage unzählige Stunden Zeit bei Struktur- und Designfragen, damit Sie Ihre wertvolle Zeit ganz auf den Inhalt konzentrieren können.

Der Studiengang Medizinische Informatik gehört sowohl zur Medizinischen Fakultät als auch zur Fakultät für Mathematik und Informatik (FMI) der Universität Leipzig,
aber es gelten nur die \href{https://www.mathcs.uni-leipzig.de/studium/studienorganisation/formulare}{Vorgaben des Prüfungsamtes der FMI},
insbesondere die Richtlinien zur Anfertigung von Bachelor- Master- und Diplomarbeiten.
Darüber hinaus ist die Vorlage an die speziellen Anforderungen der Arbeitsgruppe \href{https://www.imise.uni-leipzig.de/en/Research/MaGIS}{MaGIS} des \href{https://www.imise.uni-leipzig.de/}{IMISE} angepasst, sie kann aber auch darüber hinaus verwendet werden.
Insbesondere bestehen wir auf eine vorgegebene strukturierte Einleitung, siehe \cref{ch:introduction}.
Die MaGIS-Vorlage kann von den Hinweisen des Prüfungsamtes leicht abweichen, z.B. finden wir diese Kurzzusammenfassung am Anfang der Arbeit sinnvoller als nach dem Hauptteil, in beiden Fällen werden wir Ihnen keine Punkte abziehen.
Andere Kapitel sind jedoch flexibel, scheuen Sie sich nicht, nach Rücksprache mit ihren Betreuern bei Bedarf davon abzuweichen.
Falls Sie Fehler finden oder Änderungswünsche haben, gerne im \href{https://github.com/IMISE/imise-classicthesis/}{Git Repository} eine Issue öffnen oder ein Pull Request öffnen.

Diese Vorlage ist selbstbeschreibend:
Lesen Sie jeden Abschnitt gründlich und ersetzen Sie die Hilfestellungen durch Ihren eigenen Text.

\iffalse
In der Zusammenfassung wird vor allem zusammenfassend dargestellt, welche Ergebnisse zu den formulierten Zielen erreicht wurden.
Bei der Formulierung sollten wichtige Schlüsselbegriffe verwendet werden.
Wurden Fragen formuliert, können auch diese hier beantwortet werden.
So soll es möglich sein, dass ein Leser von der Arbeit lediglich die Einleitung und die Zusammenfassung lesen und doch die Ergebnisse der Arbeit erfassen kann.
\fi
